\chapter{Perception Interfaces Before VLA}

\section*{학습 목표}
Chapter 2는 observation, state, belief의 차이를 고정했다. 이 장은 그 사이를 채우는 perception interface를 다룬다. VLA 이전의 robot perception은 image를 곧바로 action으로 바꾸지 않았다. 중간 표현을 만들고, 그 표현을 planning, control, safety, evaluation이 소비했다. 이 장을 읽고 나면 독자는 다음을 이해해야 한다.

\begin{enumerate}
  \item camera image가 robot state가 되기까지 어떤 geometric, semantic, object-centric interface가 필요한지 설명할 수 있다.
  \item calibration, depth, pose, mask, map, uncertainty가 controller와 planner에 어떤 계약을 제공하는지 이해한다.
  \item deep perception이 명시적 perception contract를 없앤 것이 아니라 box, mask, feature map, neural field 같은 다른 표현으로 바꾼 것임을 안다.
  \item 최신 VLA가 perception responsibility를 hidden representation 안으로 흡수할 때 생기는 failure attribution, safety, evaluation 문제를 읽을 수 있다.
\end{enumerate}

\section{왜 perception 장이 필요한가}

VLA 논문은 종종 ``image와 language를 입력으로 받아 action을 출력한다''고 말한다. 모델 구조를 설명할 때는 이 문장이 충분할 수 있다. 그러나 robot system을 설명하기에는 부족하다. 로봇이 실제 세계에서 움직이려면 image는 최소한 다음 질문에 답해야 한다.

\begin{itemize}
  \item 이 pixel은 어느 camera frame에서 온 것인가?
  \item object는 어디에 있고, 어느 정도 확실한가?
  \item target object와 obstacle은 metric space에서 얼마나 떨어져 있는가?
  \item end-effector가 접근할 수 있는 pose와 grasp region은 어디인가?
  \item 지금 보이는 scene representation은 controller와 같은 시간 기준을 쓰는가?
\end{itemize}

이 질문들은 VLA 이전 로봇공학에서 perception module이 명시적으로 담당하던 책임이다. Feature matching, calibration, stereo, RGB-D fusion, SLAM, detection, segmentation, 6D pose estimation, tracking, reconstruction은 서로 다른 기술처럼 보이지만, system 관점에서는 모두 같은 질문을 푼다. 즉 image observation을 planning과 control이 사용할 수 있는 intermediate representation으로 바꾼다 \cite{hartley2004,lowe2004sift,murartal2015orbslam,ren2015fasterrcnn,he2017maskrcnn}.

\begin{definitionbox}{Perception interface}
Perception interface는 raw observation $\obs_t$를 robot stack이 소비할 수 있는 중간 표현으로 바꾸는 계약이다. 이 계약은 coordinate frame, metric scale, object identity, pose, mask, map, uncertainty, timestamp, failure flag를 포함할 수 있다.
\end{definitionbox}

Real VLA에서 중요한 점은 VLA가 이 계약을 없앴는가가 아니다. 많은 VLA는 이 계약의 일부를 model embedding이나 policy hidden state 안으로 흡수한다. 그 결과 system diagram은 단순해질 수 있지만, 책임이 사라지지는 않는다. 오히려 실패 분석은 더 어려워진다. 컵을 잘못 집었을 때 원인이 language grounding인지, object detection인지, depth error인지, camera calibration인지, controller handoff인지 분리하기 어려워진다.

\section{Perception contract의 기본 형식}

가장 단순한 form으로 쓰면 perception module은 다음 mapping을 제공한다.

\begin{equation}
  \phi_t
  =
  \mathrm{Perceive}(\obs_{1:t}, \conf_{1:t}, \mathcal{K}, \mathcal{T}),
  \label{eq:perception-map}
\end{equation}

여기서 $\mathcal{K}$는 camera intrinsics, $\mathcal{T}$는 camera, robot base, end-effector, world frame 사이의 extrinsic calibration이다. 출력 $\phi_t$는 하나의 vector가 아니라 structured state evidence이다.

\begin{equation}
  \phi_t =
  \{G_t, O_t, P_t, M_t, U_t, \tau_t\}.
  \label{eq:perception-contract}
\end{equation}

$G_t$는 geometry evidence, $O_t$는 object-level evidence, $P_t$는 pose 또는 keypoint evidence, $M_t$는 map 또는 reconstructed scene, $U_t$는 uncertainty, $\tau_t$는 timestamp와 synchronization metadata이다. 이 notation은 엄밀한 standard가 아니라 diagnostic tuple이다. 목적은 perception output을 ``feature'' 또는 ``embedding'' 하나로 뭉개지 않고, downstream module이 실제로 무엇을 받는지 분해하는 데 있다.

\begin{table}[ht]
\centering
\caption{Perception output과 downstream responsibility}
\label{tab:perception-contract}
\small
\begin{tabularx}{\textwidth}{p{25mm}X X}
\toprule
Output & 예시 & 소비하는 module \\
\midrule
Geometry $G_t$ & depth, point cloud, surface normal, free space & motion planning, collision checking, grasp synthesis \\
Object $O_t$ & class, instance id, bounding box, mask & language grounding, task planning, object selection \\
Pose $P_t$ & 6D pose, keypoints, articulated pose, grasp frame & controller target, manipulation primitive, tracking \\
Map $M_t$ & occupancy, TSDF, visual map, neural field & navigation, scene memory, long-horizon planning \\
Uncertainty $U_t$ & covariance, confidence, ambiguity set & safety margin, active perception, fallback decision \\
Time $\tau_t$ & sensor timestamp, latency, dropped frame & real-time control, logging, evaluation \\
\bottomrule
\end{tabularx}
\end{table}

VLA가 end-to-end policy로 학습되더라도 이 표의 항목은 암묵적으로 남는다. 예를 들어 action이 Cartesian delta pose라면 controller는 target frame과 metric scale을 필요로 한다. Safety shield가 forbidden region을 검사하려면 object pose와 uncertainty가 필요하다. Evaluation이 ``glass를 건드리지 않았다''고 말하려면 glass geometry와 contact 또는 proximity evidence가 필요하다.

\section{Calibration과 projection}

전통 perception의 첫 번째 계약은 image coordinate와 robot coordinate를 연결하는 것이다. Camera point $\bm{x}^C$를 robot base frame의 point $\bm{x}^B$로 바꾸려면 extrinsic transform이 필요하다.

\begin{equation}
  \bm{x}^{B} = T_{BC}\bm{x}^{C},
  \qquad
  T_{BC}\in SE(3).
  \label{eq:perception-frame-transform}
\end{equation}

Image pixel $\bm{p}=(u,v,1)^\top$와 camera ray는 intrinsics $K$를 통해 연결된다.

\begin{equation}
  \lambda \bm{p} = K \bm{x}^{C}.
  \label{eq:camera-projection}
\end{equation}

이 식은 고전적인 projective geometry처럼 보이지만, VLA에서 매우 실용적인 의미를 가진다 \cite{hartley2004}. 모델이 ``컵 오른쪽을 잡아라''는 action을 냈을 때, 그 오른쪽이 image plane의 오른쪽인지, object frame의 오른쪽인지, robot base frame의 오른쪽인지가 다르면 실제 motion은 달라진다. Calibration error가 있으면 VLA policy가 semantic target을 맞게 골라도 controller는 틀린 metric location으로 움직인다.

\begin{warningbox}{Semantic correctness는 metric correctness가 아니다}
VLM이 target object를 정확히 가리켜도, camera intrinsics, extrinsics, depth scale, timestamp가 틀리면 robot action은 실패한다. Real VLA claim은 object recognition claim과 robot execution claim을 분리해서 읽어야 한다.
\end{warningbox}

따라서 VLA paper를 읽을 때 camera setup이 단순한 실험 환경 정보로 보이면 안 된다. Camera placement, calibration procedure, hand-eye calibration, image crop, resize, normalization은 action representation과 같은 system contract의 일부이다.

\section{Depth, point cloud, free space}

Manipulation과 navigation은 metric space에서 일어난다. Depth는 image를 metric geometry로 바꾸는 핵심 interface이다. RGB-D camera, stereo, multi-view reconstruction, monocular depth estimator는 구현은 다르지만 downstream에는 비슷한 질문을 제공한다.

\begin{equation}
  D_t(u,v) \mapsto \bm{x}^{C}_{u,v}
  =
  D_t(u,v)K^{-1}
  \begin{bmatrix}
  u\\v\\1
  \end{bmatrix}.
  \label{eq:depth-unprojection}
\end{equation}

Depth map은 point cloud $\mathcal{P}_t=\{\bm{x}_i\}$로 바뀌고, point cloud는 surface, obstacle, free space, graspable region을 추정하는 데 사용된다. Control 관점에서는 이 정보가 collision constraint와 workspace constraint를 만든다.

\begin{equation}
  \mathcal{X}_{free,t}
  =
  \{\bm{x}\in\mathcal{X}\mid d(\bm{x},\mathcal{P}_t) \ge d_{\min}\}.
  \label{eq:free-space}
\end{equation}

여기서 $d(\bm{x},\mathcal{P}_t)$는 point cloud 또는 reconstructed surface까지의 거리이다. VLA policy가 waypoint를 제안하더라도, motion planner와 safety layer는 \cref{eq:free-space} 같은 geometric free-space contract를 필요로 한다.

최신 VLA에서는 depth가 명시적으로 주어지지 않는 경우가 많다. RGB image와 language만으로 action을 생성할 수 있다면 system은 단순해진다. 그러나 depth responsibility가 사라진 것은 아니다. Hidden representation 안에 metric relation이 충분히 들어 있는지, 아니면 training distribution의 scene layout shortcut을 쓰는지 평가해야 한다.

\section{Detection, segmentation, tracking}

Object-centric manipulation에서는 ``무엇을 움직일 것인가''와 ``어디를 피할 것인가''가 중요하다. 전통적인 perception stack은 이를 detection, segmentation, tracking으로 분리했다. Detection은 object 후보와 class를 만들고, segmentation은 pixel-level support를 제공하며, tracking은 time consistency를 유지한다 \cite{ren2015fasterrcnn,he2017maskrcnn}.

\begin{equation}
  O_t =
  \{(c_i, B_i, S_i, \gamma_i)\}_{i=1}^{N_t},
  \label{eq:object-set}
\end{equation}

여기서 $c_i$는 class 또는 semantic label, $B_i$는 bounding box, $S_i$는 mask, $\gamma_i$는 confidence이다. 이 object set은 language grounding의 입력이 된다.

\begin{equation}
  i^\star
  =
  \arg\max_i
  \mathrm{Ground}(\task, c_i, S_i, h_t).
  \label{eq:object-grounding}
\end{equation}

VLM/VLA는 \cref{eq:object-grounding}을 명시적으로 수행하지 않을 수 있다. Image token과 text token의 cross-attention 안에서 target object를 implicit하게 고를 수 있다. 하지만 Real VLA evaluation에서는 여전히 object selection error를 분리해야 한다. Mug/rack/glass 예제에서 실패는 적어도 다음으로 나뉜다.

\begin{itemize}
  \item mug를 glass로 잘못 grounding한 failure
  \item mug mask는 맞지만 depth 또는 pose가 틀린 failure
  \item mug pose는 맞지만 grasp affordance가 틀린 failure
  \item action은 맞지만 controller가 tracking하지 못한 failure
\end{itemize}

명시적 detection/segmentation pipeline은 이런 failure attribution을 쉽게 만든다. End-to-end VLA는 이런 attribution을 추가 log, attention analysis, intervention label, object-level probe로 다시 확보해야 한다.

\section{Pose, keypoint, grasp frame}

Robot control은 object label만으로 움직일 수 없다. Manipulation에는 pose 또는 grasp frame이 필요하다. Rigid object에 대해서는 6D pose를 다음처럼 쓸 수 있다.

\begin{equation}
  P_t^i = (R_t^i,\bm{p}_t^i),\qquad R_t^i\in SO(3),\;\bm{p}_t^i\in\mathbb{R}^3.
  \label{eq:object-pose}
\end{equation}

Grasp candidate는 object pose와 local grasp library 또는 learned grasp head로부터 나온다.

\begin{equation}
  g_t^j = (T_{BG}^j, w_j, \eta_j),
  \label{eq:grasp-candidate}
\end{equation}

여기서 $T_{BG}^j$는 base frame에서 본 grasp pose, $w_j$는 gripper width, $\eta_j$는 grasp quality score이다. VLA가 ``pick mug''라는 high-level action을 낸다면, 실제 executor는 \cref{eq:grasp-candidate} 같은 grasp-level interface를 필요로 할 수 있다.

Keypoint representation도 중요하다. Drawer handle, tool tip, deformable object, articulated object에서는 full object pose보다 task-relevant keypoint가 더 직접적일 수 있다.

\begin{equation}
  K_t = \{\bm{k}_{1,t},\ldots,\bm{k}_{m,t}\},
  \qquad
  \act_t = f_{\theta}(K_t,\task,\conf_t).
  \label{eq:keypoint-action}
\end{equation}

VLA가 keypoint를 명시적으로 출력하지 않아도, policy hidden state는 이런 task-relevant geometry를 내부적으로 가져야 한다. 문제는 그 geometry가 system 밖으로 노출되지 않으면 controller handoff와 safety review가 약해진다는 점이다.

\section{Mapping and reconstruction}

Long-horizon task에서는 현재 frame 하나로 충분하지 않다. Robot은 이전에 본 object, obstacle, free space, target region을 기억해야 한다. Traditional robotics는 이를 map으로 다루었다. Visual SLAM은 camera trajectory와 sparse/dense map을 함께 추정한다 \cite{murartal2015orbslam}.

\begin{equation}
  M_t =
  \mathrm{MapUpdate}(M_{t-1}, \obs_t, \hat{T}_{WB,t}).
  \label{eq:map-update}
\end{equation}

Mapping representation은 여러 형태를 가진다.

\begin{itemize}
  \item sparse feature map: localization과 camera tracking에 유리하다.
  \item occupancy grid 또는 voxel map: collision checking과 navigation에 유리하다.
  \item TSDF 또는 mesh: surface reconstruction과 manipulation planning에 유리하다.
  \item neural field 또는 radiance field: view synthesis와 dense scene memory에 유리하다 \cite{mildenhall2020nerf}.
\end{itemize}

VLA에서는 scene memory가 language-conditioned hidden state, retrieval memory, image history, external map으로 나뉠 수 있다. Real VLA 관점에서 중요한 질문은 map이 어떤 형태인가보다, 그 map이 downstream claim을 닫는가이다. 예를 들어 ``부엌에서 컵을 찾아 drying rack에 놓아라''는 task는 object permanence와 workspace memory를 요구한다. 이때 map이 없다면 model은 현재 frame 밖의 object와 obstacle을 어떻게 다루는지 설명해야 한다.

\section{Affordance and contact-relevant perception}

Perception이 object를 찾는 것에서 끝나면 manipulation은 약하다. Robot은 object가 어떻게 잡히고, 밀리고, 열리고, 변형되는지 알아야 한다. 이를 affordance 또는 contact-relevant perception으로 볼 수 있다.

\begin{equation}
  A_t(\bm{x}) =
  p(\text{success}\mid \bm{x}, \obs_t, \task, \conf_t),
  \label{eq:affordance-field}
\end{equation}

$A_t(\bm{x})$는 특정 위치 또는 pose에서 action이 성공할 가능성을 나타내는 affordance field로 읽을 수 있다. SayCan류 구조에서는 language feasibility와 affordance를 결합하고, CLIPort나 Transporter 계열은 language-conditioned pick/place heatmap을 만든다 \cite{saycan2022,cliport2022,florence2021transporter}. VLA policy도 유사한 정보를 action head 안에 흡수할 수 있다.

그러나 affordance는 semantic label과 다르다. ``handle''이라는 label을 아는 것과 drawer를 여는 force direction, grasp width, contact compliance를 아는 것은 다르다. Contact-rich task에서는 perception output이 다음까지 이어져야 한다.

\begin{itemize}
  \item contact surface와 normal
  \item expected friction 또는 slip risk
  \item grasp closure 가능성
  \item articulated joint axis
  \item deformability 또는 compliance
\end{itemize}

이 정보가 없으면 VLA는 plausible action을 내지만 실행은 실패한다. Chapter 5의 contact와 force 장은 여기서 넘어온 perception evidence를 사용한다.

\section{Uncertainty and active perception}

전통 perception stack의 큰 장점은 uncertainty를 외부로 낼 수 있다는 점이다. Object pose estimate를 평균만 내는 것이 아니라 covariance와 함께 낼 수 있다.

\begin{equation}
  P_t^i \sim \mathcal{N}(\hat{P}_t^i, \Sigma_{P,t}^i).
  \label{eq:pose-uncertainty}
\end{equation}

Safety margin과 planning constraint는 uncertainty에 따라 달라져야 한다.

\begin{equation}
  \mathcal{C}_{safe,t}
  =
  \{\state\mid
  d(\state,\hat{O}_t) - \lambda \sigma_d(\state) \ge d_{\min}
  \}.
  \label{eq:perception-safety-constraint}
\end{equation}

Uncertainty가 크면 robot은 바로 실행하지 않고 active perception을 선택할 수 있다.

\begin{equation}
  \act_t \in
  \{\text{move-camera},\text{ask-human},\text{relocalize},\text{execute}\}.
  \label{eq:active-perception-action}
\end{equation}

이 식은 VLA에서 특히 중요하다. Language-conditioned policy는 uncertainty가 높은 상황에서 ``더 보기'', ``다시 확인하기'', ``실행을 멈추기'' 같은 action을 낼 수 있어야 한다. Benchmark가 success trajectory만 측정하면 이런 ability는 잘 보이지 않는다. Real VLA evaluation은 uncertainty-aware behavior를 별도 evidence로 기록해야 한다.

\section{From explicit perception to hidden representation}

이제 핵심 질문으로 돌아가자. VLA는 explicit perception module을 대체하는가? 일부 task에서는 그렇다. Large-scale image-language pretraining과 robot demonstration은 object identity, spatial relation, affordance, action prior를 하나의 representation 안에 섞을 수 있다. 이 구조는 engineering overhead를 줄이고, 기존 pipeline이 놓치는 open-vocabulary semantics를 더 잘 처리할 수 있다 \cite{clip2021,flamingo2022,blip22023,openvla2024}.

하지만 Real VLA 관점에서는 다음 trade-off를 반드시 봐야 한다.

\begin{table}[ht]
\centering
\caption{명시적 perception interface와 hidden VLA representation의 trade-off}
\label{tab:explicit-hidden-perception}
\small
\begin{tabularx}{\textwidth}{p{28mm}X X}
\toprule
항목 & 명시적 perception interface & hidden VLA representation \\
\midrule
Open vocabulary & 사전 정의된 detector에 제한될 수 있음 & language prior와 web-scale visual prior 활용 가능 \\
Metric geometry & calibration, depth, pose가 명시됨 & 내부 표현에 흡수되어 검증이 어려움 \\
Failure attribution & detection, pose, tracking, control failure 분리 쉬움 & end-to-end failure로 뭉개질 위험 \\
Safety check & geometry와 uncertainty를 shield가 직접 사용 가능 & safety layer가 별도 evidence를 요구함 \\
Adaptation & module별 교체 가능 & data와 fine-tuning에 의존 \\
Evaluation & intermediate metric 평가 가능 & probing, logging, intervention label 필요 \\
\bottomrule
\end{tabularx}
\end{table}

따라서 이 책의 입장은 pipeline으로 돌아가자는 것이 아니다. VLA가 perception을 흡수할 수 있다면 좋은 일이다. 다만 흡수된 responsibility는 없어지지 않는다. 그것은 evaluation protocol, safety monitor, dataset metadata, failure log에서 다시 드러나야 한다.

\section{Running example: mug, rack, glass}

``머그컵을 drying rack에 올리고 유리컵은 건드리지 말라''는 task를 perception interface 관점에서 다시 쓰면 다음과 같다.

\begin{equation}
  \phi_t =
  \{
  O_t^{mug},
  O_t^{rack},
  O_t^{glass},
  P_t^{mug},
  M_t^{free},
  U_t,
  \tau_t
  \}.
  \label{eq:mug-perception-contract}
\end{equation}

여기서 필요한 perception evidence는 단순히 ``mug가 보인다''가 아니다.

\begin{itemize}
  \item mug instance와 glass instance가 구분되어야 한다.
  \item mug의 graspable region과 pose가 충분히 정확해야 한다.
  \item rack slot 또는 place region이 free space로 확인되어야 한다.
  \item glass와 forbidden region의 위치 및 uncertainty가 safety layer로 전달되어야 한다.
  \item camera frame과 robot base frame 사이 transform이 controller handoff에 맞아야 한다.
  \item execution 중 object tracking 또는 re-detection이 실패하면 recovery 또는 stop이 가능해야 한다.
\end{itemize}

VLA가 이 모든 것을 hidden state 안에서 처리할 수도 있다. 그러나 reporting에서는 적어도 다음 질문을 해야 한다.

\begin{enumerate}
  \item 실패가 object selection failure였는가?
  \item 실패가 pose/depth/calibration failure였는가?
  \item 실패가 action representation 또는 controller failure였는가?
  \item glass near miss가 perception uncertainty와 함께 기록되었는가?
\end{enumerate}

이 질문이 없으면 성공률은 높아도 Real VLA claim은 약하다.

\section{Perception checklist for VLA papers}

VLA 논문을 읽을 때 perception 관련 claim은 다음 체크리스트로 확인한다.

\begin{enumerate}
  \item Camera setup, intrinsics, extrinsics, calibration procedure가 명시되는가?
  \item RGB만 쓰는지, depth/RGB-D/point cloud/proprioception을 함께 쓰는지 분명한가?
  \item Object identity, mask, pose, map, affordance 중 무엇이 explicit하고 무엇이 implicit한가?
  \item Action이 어느 coordinate frame에서 정의되는가?
  \item Hidden representation이 metric geometry를 학습했다는 evidence가 있는가?
  \item Object selection failure와 control failure를 분리해서 보고하는가?
  \item Depth, pose, occlusion, lighting, camera shift에 대한 stress test가 있는가?
  \item Uncertainty, confidence, ambiguity가 safety decision에 들어가는가?
  \item Perception failure가 dataset label 또는 intervention log로 남는가?
  \item Sim benchmark와 real camera deployment 사이의 calibration gap을 설명하는가?
\end{enumerate}

\section{요약}

전통 robot perception은 image를 action으로 바로 바꾸지 않았다. It built contracts. Calibration은 pixel과 robot frame을 연결했고, depth와 reconstruction은 metric geometry를 만들었으며, detection과 segmentation은 object-centric grounding을 제공했다. Pose, keypoint, grasp frame은 controller target이 되었고, map과 uncertainty는 planning과 safety의 입력이 되었다.

최신 VLA는 이 중간 표현의 일부를 VLM/VLA representation 안으로 흡수한다. 이것은 강력한 방향이다. 그러나 Real VLA에서 흡수된 책임은 사라지지 않는다. 명시적 perception module이 없어질수록 evaluation과 logging은 더 강해야 한다. 다음 장의 control interface는 이 perception evidence가 policy action과 controller command 사이에서 어떻게 물리 명령으로 바뀌는지를 다룬다.

\chapterwork
{Hartley and Zisserman의 camera geometry, ORB-SLAM의 metric scene reconstruction, Faster R-CNN/Mask R-CNN의 object interface, CLIPort/Transporter의 language-conditioned spatial action interface를 함께 읽어라.}
{VLA가 RGB image만 입력으로 사용한다고 할 때, 그 system이 metric geometry, object identity, uncertainty를 어디에 저장하거나 어떻게 검증하는가?}
{mug/rack/glass task에 대해 explicit perception contract $\phi_t=\{G_t,O_t,P_t,M_t,U_t,\tau_t\}$를 설계하고, 각 항목이 planner, controller, safety layer 중 어디로 전달되는지 표로 정리하라.}
